\chapter{Estado del arte}
\label{cap:estado-arte}

Este capítulo revisa los trabajos previos relevantes en dos líneas
complementarias: el análisis automático de mamografías mediante aprendizaje
profundo, y las aplicaciones del Análisis Topológico de Datos en imagen médica.

\section{Análisis automático de mamografías}
El análisis asistido por computador de mamografías tiene una larga trayectoria,
desde los primeros sistemas CAD basados en características diseñadas manualmente
(\textit{hand-crafted features}) hasta los actuales enfoques de aprendizaje
profundo. Las redes neuronales convolucionales (CNN) se han consolidado como el
paradigma dominante en visión por computador y en imagen médica, gracias a su
capacidad para aprender jerarquías de características directamente de los
datos~\cite{litjens2017survey, alieli2024survey}.

En el ámbito específico de la mamografía, destacan trabajos orientados a la
detección y clasificación de masas y microcalcificaciones, así como a la
estimación del riesgo. Yala \textit{et al.}~\cite{yala2019} propusieron un
modelo basado en \textit{deep learning} que mejora la predicción del riesgo de
cáncer de mama a partir de la mamografía, superando a los modelos de riesgo
tradicionales. El uso de \textit{transfer learning} a partir de redes
preentrenadas en ImageNet (ResNet, DenseNet, EfficientNet, etc.) es una práctica
habitual para mitigar la escasez de datos etiquetados en el dominio médico.

\subsection{Conjuntos de datos públicos}
Entre los conjuntos de datos públicos de mamografía más utilizados destaca
\mbox{CBIS-DDSM}~\cite{lee2017cbisddsm}, una versión curada y estandarizada del
clásico DDSM que incluye segmentaciones de ROI verificadas y etiquetas de
patología. Otros recursos relevantes son INbreast y las bases de datos de
ecografía mamaria.

\subsection{Limitaciones}
A pesar de sus buenos resultados, los modelos de \textit{deep learning} en
mamografía presentan limitaciones bien conocidas:
\begin{itemize}
  \item \textbf{Interpretabilidad limitada}: su naturaleza de caja negra
        dificulta la confianza clínica y la certificación.
  \item \textbf{Dependencia de grandes volúmenes de datos} etiquetados.
  \item \textbf{Riesgo de fuga de datos} (\textit{data leakage}) y de sesgos
        de adquisición, que pueden producir métricas artificialmente altas si el
        protocolo experimental no es riguroso.
\end{itemize}
Estas limitaciones motivan la búsqueda de representaciones más robustas e
interpretables, como las que proporciona el TDA.

\section{Análisis Topológico de Datos en imagen médica}
El Análisis Topológico de Datos (TDA) engloba un conjunto de técnicas que
estudian la forma de los datos mediante herramientas de topología
algebraica~\cite{carlsson2009topology, edelsbrunner2010computational}. Su
herramienta central, la \emph{homología persistente}, ha sido aplicada con éxito
en diversos problemas de imagen médica: caracterización de texturas de tumores,
análisis de imágenes histopatológicas, neuroimagen y detección de estructuras en
radiología.

En el contexto de la mamografía y del cáncer de mama, los descriptores
topológicos permiten capturar patrones de conectividad y estructura (número de
componentes conexas, agujeros, ramificaciones de conductos, distribución de
microcalcificaciones) que resultan difíciles de modelar con características
convencionales. Su robustez frente al ruido y su interpretabilidad geométrica
los convierten en un complemento natural de las CNN.

\section{Síntesis y aportación de este trabajo}
La revisión anterior muestra que, si bien el \textit{deep learning} domina el
análisis de mamografías, el TDA ofrece una vía prometedora y poco explorada para
obtener representaciones interpretables y robustas. La aportación de este TFG es
un estudio sistemático de descriptores topológicos sobre \mbox{CBIS-DDSM}, su
evaluación como características para la clasificación benigno/maligno, y su
comparación y combinación con un modelo de \textit{deep learning} de referencia.
